% Page 055

\seite{55}

\abschnitt{Schuljahr 1914–18}
\pstart
\margmark{Mit dem 18. Feb\\2 Knaben und zw}%
ruar wurden 4 Kinder aus der Schule entlassen\ob
ei Mädchen, welche sich alle der Landwirtschaft

Kriegsanleihe\ob
\margmark{die hiesige Sch\\"        "\\"        "}%
ule zeichnete zur fünften Kriegsanleihe 352 M.\ob
 "     "    6.     "                    720 M.\ob
 "     "    7.     "                    580 "

Schuljahr 1917–18\ob
\margmark{Es wurden 9 Kin\\4 Knaben, 5 Mäd\\sind 14 Knaben,}%
der in die Schule aufgenommen\ob
chen. Die Schülerzahl beträgt 38, davon\ob
 24 Mädchen.

\pend
\abschnitt{Ernte}
\pstart

\margmark{Die Gesamternte\\allgemeinen sei\\beiteten. Heu g\\mittelte Leute\\Früchte als Ers\\waren sehr hoch}%
 war befriedigend. Aber die Getreideernte im\ob
 deshalb besser weil jetzt alle mehr ar\ohb
ab's stellenweise nur wenig, so daß selbst be\ohb
kaum für ihren Bedarf hatten, und vielerlei\ob
atzfrüchte anwenden mußten. Die Viehpreise\ob
. Der Preis der Lebensmittel stieg fortwährend.

\pend
\abschnitt{Schuljahr 1918–19}
\pstart

\margmark{Entlassen wurde\\3 Mädchen und e\\bei der Landwir\\Ostern 17.4.19 3 K\\gen, 13 Knaben}%
n in diesem Jahre 4 Kinder.\ob
in Knabe, sämtliche halfen\ob
tschaft. Aufgenommen wurden\ob
inder. die Schülerzahl ist mit 39 Schulpflichti\ohb
und 26 Mädchen.

\pend
\abschnitt{Ernte}
\pstart

\margmark{Die Gesamternte\\einen guten Ert\\gut, auch die V\\fortwährend.}%
 war eine mittlere Ernte. Dem Umfange nach\ob
rag. Hafer, altdeutsche Sommerfrucht, gut, sehr wenig Kartoffeln\ob
iehpreise waren hoch. Der Preis der Lebensmittel stieg\ob
\pend
